\chapter{Conclusion and outlook} \label{chap:summary}

In this final chapter the project under study and its main findings are summarised.
Then, the limitations of the study are outlined and directions for future work are suggested. 

In this thesis, \acrlong{tbl}s (TBL) along curved walls were investigated.\fxchange{Past  tense, für alles as gemacht wurde}
The studies were performed using simulations based on \acrfull{dns} and \acrfull{rans}.
Based on these simulations, an analysis of the impact of convex curvatures on the development of the \acrlong{tbl} was provided.
However, in flows along a curvature, a pressure gradient naturally arises, induced by the geometry.\fxchange{Fakten und Resultate: Present}
The pressure gradient also affects the \acrlong{tbl}.
In order to study the effect of the curvature separately,  it is necessary to know the effect of the pressure gradient.
Consequently, a manipulation of the pressure field is required.
In our approach we manipulated the pressure such that a \acrfull{zpg} develops.
This is a favourable pressure distribution, since \acrlong{tbl} along flat plate walls under \acrlong{zpg} conditions are well-studied and provide an established reference setup.


The generation of a \acrshort{zpg} in a curved simulation domain was achieved through an optimisation process.
We adopted the optimisation algorithm introduced by \citet{morita_applying_2022} that performs a shape optimisation in a \acrshort{rans} simulation.
\fxunsure{Wenn hier das etablierte Vorgehen nach Morita2022 erklärt wird, dann Persent lassen, wenn das jetzt Dein Vorgehen ist, d.h. das ist eine Anpassung auf Curvatures, dann "For our study," und alles in Past}
The upper wall is characterised by multiple  \benennung{control points}, which are interconnected by a spline.
The position of these \benennung{control points} is optimised such that the Clauser parameter $\beta$ at the lower wall fits its target value.
The Clauser parameter corresponds to the non-dimensionalisation of the pressure gradient.
Through the adaptation of the \benennung{control points} the shape of the domain is modified, which implicitly affects the pressure distribution.
Once the optimisation has converged, the result s applied on a domain for a \acrshort{dns}.
In order to accomplish this objective, the velocities along a designated wall-normal height are extracted.
The height corresponds to the height of the upper wall of the domain used for \acrshort{dns}.
In this context, the extracted velocities are imposed as  \benennung{Dirichlet} \acrlong{bc}.
%
\fxunsure{Ab hier ist es in jedem Fall Dein Ansatz: Daher jetzt definitiv: Past}
The optimised \acrlong{bc} were employed to simulate three distinct  curved setups.
Two setups are characterised by a radius of $100\delta_0$ with curvatures of $30$° and $90$°.
The third setup exhibits a radius of $500 \delta_0$ and an $18$° curvature.
All setups are characterised by the presence of a straight inflow region prior to the onset of curvature. This ensures the proper\fxinfo{adequate} development of the \acrshort{tbl}.\fxinfo{das ist mE wieder Faktenbeschreibung}
The \acrshort{tbl} was generated with an induced tripping of a laminar Blasius boundary layer in the proximity of the inflow.
At the termination of the curvature, a straight outlet region was  incorporated in the tangential direction, serving as a rebalancing zone.
 Two simulations were performed for each curved case using these setups. 
 As a base configuration, all cases were  simulated with a \benennung{Dirichlet} \acrlong{bc} of $(1,0,0)$ velocity profile at the upper wall.
Then,  simulations were performed with the optimised velocity \acrlong{bc} at the upper wall.
With the base simulations as a reference, the effect of the optimisation could be studied.
In addition, the optimisation capability was compared for the three curvature cases.

\fxinfo{Resultate: present nutzen}
The results of our study show that the optimisation capability differs for the \benennung{local} and \benennung{global} trends of the pressure distribution.
On a \benennung{global} level, the optimised versions for all three cases demonstrate a more constant\fximprove{uniform/stable?} pressure distribution in the curvature when compared to their respective base cases.
Nonetheless, fluctuations in the pressure distribution manifest  \benennung{locally}, and are also present in the distribution of the Clauser parameter.
The effects of different curvatures are studied through a comparison of case R100deg90 and R500deg18.
XXX curvature effects 
XXX optimisation for different cases


%In this thesis, the workflow, introduced in \citet{schlatter_leveraging_2025}, to generate a curved \acrshort{tbl} domain with \acrshort{zpg} was tested.
%\fxunsure{Jetzt bin ich verwirrt: das hatten wir doch gerade schon..., ich lösch den Satz}
The findings of this study provide evidence that the application of the optimisation algorithm presented by \citet{morita_applying_2022} is feasible across all three curvature configurations.
The \benennung{global} trend of the pressure field was successfully manipulated towards a constant value in the curved segment.
However, \benennung{local} fluctuations in the pressure distribution were detected.
The positioning of the fluctuations suggests a correlation with the location  of the \benennung{control points} used during  in the optimisation.
The studies conducted by \citet{ganagdhar_bayesian_2025} examined the optimisation capability and identified a case characterised by a substantially diminished curvature in comparison to the boundary layer thickness. \fxinfo{wieso ist das hier irgendwie ein (nicht) vergleichbarer Fall?}
Therefore, to our knowledge there is only limited knowledge  on the capability of the optimisation algorithm in relation to the geometric properties examined here.
In order to eliminate the \benennung{local} fluctuations in the pressure distribution, a detailed study of the optimisation algorithm with various settings and geometries is required.
Furthermore, an investigation into the number and positioning of the \benennung{control points} is necessary in order to establish an optimal relation between a smooth $\beta$ distribution and a reasonable computational effort for the optimisation.


\fxunsure{Geht es hier um andere "fluctuations " als die local  fluctuations die schon besprochen wurden? oder ist das eher eine Wiederholung?}
In the \acrshort{dns} of the curved \acrshort{tbl},  fluctuations in between the instantaneous pressure fields were detected.
Due to their characteristics analysed in \autoref{sec:pressFluc}, we hypothesised  that these are non-physical artefacts caused by the outer \fxinfo{external?} constraints of the domain.
Similar pressure modes were observed by \citet{maday_spectral_1989} and \citet{sabbah_filtering_2001} who introduced methods to eliminate them.
We performed  a series of tests with diverse configurations with the objective of identifying the underlying causes of the observed fluctuations.
However, the conducted tests did not allow for the clear identification of  the causes of the fluctuations.
In order to determine the root cause of these fluctuations, further tests are required.
Nevertheless, we verified the reliability of the statistical results to ensure that the statistics are unaffected by the fluctuations.
In order to accomplish this objective,   we extracted data from the straight section and compared them with the earlier performed simulations of a plane \acrshort{tbl} as well as with reference data from plane \acrshort{tbl} simulations.
This verification was presented in \autoref{sec:planeTBLVerif}. It proves  that the statistical quantities are undisturbed and hence similar to the reference data.
In order to obtain a  more profound  of the fluctuations, further tests are required in future research.
This will also help for future enhancement to  the domain setups, as well as facilitating the comparison of different setups.

\fxunsure{Soll die Bibliographie so aussehen? Bitte nochmal draufschauen}
